
\documentclass{article}
\usepackage{colortbl}
\usepackage{makecell}
\usepackage{multirow}
\usepackage{supertabular}

\begin{document}

\newcounter{utterance}

\centering \large Interaction Transcript for game `cladder', experiment `full\_v1.5\_default', episode 6687 with Qwen3.6{-}27B.
\vspace{24pt}

{ \footnotesize  \setcounter{utterance}{1}
\setlength{\tabcolsep}{0pt}
\begin{supertabular}{c@{$\;$}|p{.15\linewidth}@{}p{.15\linewidth}p{.15\linewidth}p{.15\linewidth}p{.15\linewidth}p{.15\linewidth}}
    \# & \multicolumn{2}{c}{Player} && \multicolumn{2}{c}{Game Master} \\
    \hline

    \theutterance \stepcounter{utterance}  
    & & & \multicolumn{4}{p{0.6\linewidth}}{
        \cellcolor[rgb]{0.9,0.9,0.9}{
            \makecell[{{p{\linewidth}}}]{
                \texttt{\tiny{[P1$\langle$GM]}}
                \texttt{You are an expert in causal inference. The following question is not a typical commonsense query, but rather a meticulously designed question created by a professor specializing in causal inference, intended to assess the students' mastery of the course content.} \\
\texttt{Imagine a self{-}contained, hypothetical world with only the following conditions, and without any unmentioned factors or causal relationships: Xevo has a direct effect on tijw and tijv. Tijv has a direct effect on gyzp. Tijw has a direct effect on gyzp. The overall probability of xevo is 48\%. For those who are not xevo, the probability of gyzp is 39\%. For those who are xevo, the probability of gyzp is 78\%. Is gyzp more likely than not gyzp overall?} \\
\texttt{Start your answer with 'Yes' or 'No', followed by additional reasoning or evidence to support your explanation.} \\
            }
        }
    }
     \\ \\

    \theutterance \stepcounter{utterance}  
    & \multicolumn{4}{p{0.6\linewidth}}{
        \cellcolor[rgb]{0.561,0.737,0.561}{
            \makecell[{{p{\linewidth}}}]{
                \texttt{\tiny{[P1$\rangle$GM]}}
                \texttt{Yes} \\
\\ 
\texttt{To determine whether `gyzp` is more likely than not `gyzp` overall, we need to calculate the total probability of `gyzp`, denoted as \$P(gyzp)\$.} \\
\\ 
\texttt{We are given the following information:} \\
\texttt{1.  \$P(xevo) = 0.48\$} \\
\texttt{2.  \$P(\textbackslash{}neg xevo) = 1 {-} P(xevo) = 1 {-} 0.48 = 0.52\$} \\
\texttt{3.  \$P(gyzp | \textbackslash{}neg xevo) = 0.39\$} \\
\texttt{4.  \$P(gyzp | xevo) = 0.78\$} \\
\\ 
\texttt{The problem states that the world is self{-}contained with only the mentioned causal relationships. This implies that `xevo` is the only variable determining the probability of `gyzp` in the context of the provided conditional probabilities. We can use the Law of Total Probability to find the marginal probability of `gyzp`:} \\
\\ 
\texttt{\$\$P(gyzp) = P(gyzp | xevo) \textbackslash{}cdot P(xevo) + P(gyzp | \textbackslash{}neg xevo) \textbackslash{}cdot P(\textbackslash{}neg xevo)\$\$} \\
\\ 
\texttt{Substituting the given values:} \\
\\ 
\texttt{\$\$P(gyzp) = (0.78 \textbackslash{}cdot 0.48) + (0.39 \textbackslash{}cdot 0.52)\$\$} \\
\\ 
\texttt{Calculating the first term:} \\
            }
        }
    }
    & & \\ \\

    \theutterance \stepcounter{utterance}  
    & & & \multicolumn{2}{p{0.3\linewidth}}{
        \cellcolor[rgb]{0.9,0.9,0.9}{
            \makecell[{{p{\linewidth}}}]{
                \texttt{\tiny{[GM$|$GM]}}
                \texttt{Yes} \\
            }
        }
    }
    & & \\ \\

    \theutterance \stepcounter{utterance}  
    & & & \multicolumn{2}{p{0.3\linewidth}}{
        \cellcolor[rgb]{0.9,0.9,0.9}{
            \makecell[{{p{\linewidth}}}]{
                \texttt{\tiny{[GM$|$GM]}}
                \texttt{yes} \\
            }
        }
    }
    & & \\ \\

    \theutterance \stepcounter{utterance}  
    & & & \multicolumn{2}{p{0.3\linewidth}}{
        \cellcolor[rgb]{0.9,0.9,0.9}{
            \makecell[{{p{\linewidth}}}]{
                \texttt{\tiny{[GM$|$GM]}}
                \texttt{game\_result = WIN} \\
            }
        }
    }
    & & \\ \\

\end{supertabular}
}

\end{document}
